% =====================================================================
%  Section 8.4: Godunov's method.  Optional: remove the \input line in
%  main.tex to drop this subsection (nothing else depends on it).
% =====================================================================

\subsection{A numerical scheme: Godunov's method}
\label{subsec:godunov}

This extension outlines Godunov's method \cite{godunov1959}, a numerical
scheme for the conservation law
$u_t+f(u)_x=0$ that uses the exact Riemann solution of
Theorem~\ref{thm:riemann-solution} as its building block, and states its
convergence to the entropy solution.

\subsubsection{Finite volume discretisation}
\label{subsubsec:finite-volume}

Divide $\mathbb{R}$ into cells $[x_{j-1/2}, x_{j+1/2}]$ of width $\Delta
x$, with cell averages
\begin{equation}
  U_j^n \approx \frac{1}{\Delta x}\int_{x_{j-1/2}}^{x_{j+1/2}} u(x, t^n)\, dx,
\end{equation}
at times $t^n = n\Delta t$. Integrating $u_t+f(u)_x=0$ over a cell and
over $[t^n,t^{n+1}]$ gives the exact identity
\begin{equation}
  U_j^{n+1} = U_j^n - \frac{\Delta t}{\Delta x}
  \left( \widehat{f}_{j+1/2} - \widehat{f}_{j-1/2} \right),
\end{equation}
where $\widehat f_{j+1/2}$ is the (time-averaged) flux through the
interface $x_{j+1/2}$, which must be approximated by a \emph{numerical
flux function} $\widehat f_{j+1/2} = \widehat f(U_j^n, U_{j+1}^n)$.

\subsubsection{The Godunov flux}
\label{subsubsec:godunov-flux}

Godunov's method defines $\widehat f$ by solving, at each interface, the
Riemann problem with left and right states $U_j^n$ and $U_{j+1}^n$
\emph{exactly}, using Theorem~\ref{thm:riemann-solution}, and evaluating
the resulting self-similar solution along $x/t=0$:
\begin{equation}
  \label{eq:godunov-flux}
  \widehat f_{j+1/2} = f\big(u^R(0;\, U_j^n, U_{j+1}^n)\big),
\end{equation}
where $u^R(\xi; u_\ell, u_r)$ denotes the value at self-similar
coordinate $\xi=x/t$ of the entropy solution from
Theorem~\ref{thm:riemann-solution} with data $(u_\ell,u_r)$. Concretely,
by \eqref{eq:riemann-shock}--\eqref{eq:riemann-rarefaction},
\begin{equation}
  u^R(0; u_\ell,u_r) =
  \begin{cases}
    u_\ell, & u_\ell>u_r,\ s=\tfrac{u_\ell+u_r}{2}>0, \\
    u_r, & u_\ell>u_r,\ s<0, \\
    u_\ell, & u_\ell\le u_r,\ u_\ell\ge0, \\
    u_r, & u_\ell\le u_r,\ u_r\le0, \\
    0, & u_\ell < 0 < u_r,
  \end{cases}
\end{equation}
the last case corresponding to $\xi=0$ lying inside the rarefaction fan of
\eqref{eq:riemann-rarefaction}, i.e. $u^R(0)=(f')^{-1}(0)=0$ for Burgers'
flux. (In the shock case with $s=0$, i.e.\ $u_\ell=-u_r$, the value of
$u^R(0)$ is ambiguous, but $f(u_\ell)=f(u_r)$ so the flux
\eqref{eq:godunov-flux} is the same either way.)

\subsubsection{Stability and convergence}
\label{subsubsec:godunov-convergence}

\begin{proposition}[CFL condition]
\label{prop:cfl}
The Godunov scheme is stable provided the time step satisfies
\begin{equation}
  \label{eq:cfl}
  \Delta t \le \frac{\Delta x}{\max_j |f'(U_j^n)|}.
\end{equation}
\end{proposition}

Condition \eqref{eq:cfl} (the Courant--Friedrichs--Lewy condition)
ensures that information does not propagate across more than one cell
per time step, consistent with the finite propagation speed of
characteristics established in Section~\ref{sec:characteristics}. Two
distinct requirements are worth separating (see \cite[Ch.~4]{leveque2002}).
The flux \eqref{eq:godunov-flux} is the exact flux through $x_{j+1/2}$ as
long as the state at that interface stays constant during the step, i.e.\
no wave from a neighbouring interface reaches it; this holds for CFL
number $\le1$, which is \eqref{eq:cfl}, and is all that stability and
monotonicity of the scheme need. The ``reconstruct--evolve--average''
picture, in which the exact solution in each cell is the superposition of
the Riemann solutions from its two interfaces, requires in addition that
these two sets of waves do not interact inside the cell, i.e.\ CFL number
$\le\tfrac12$. The value $0.45$ used below satisfies both.

% TODO: confirm the chapter/section numbers in LeVeque (Ch.~4 for Godunov's
% method, Ch.~12 for convergence of monotone schemes) and the theorem number
% in Crandall--Majda.  Note that this theorem asserts convergence, NOT a rate.
\begin{theorem}[{Convergence; \cite{crandall1980}, \cite[Ch.~12]{leveque2002}}]
\label{thm:godunov-convergence}
Under \eqref{eq:cfl}, the Godunov approximations converge, as $\Delta x,
\Delta t \to 0$, to the unique entropy solution of
Theorem~\ref{thm:oleinik-uniqueness}.
\end{theorem}

\begin{remark}
\label{rem:godunov-check}
This gives a second, independent route to the solutions constructed by
hand in Sections~\ref{sec:nonuniqueness}--\ref{sec:riemann-problem}: an
implementation of \eqref{eq:godunov-flux} applied to the Riemann data
\eqref{eq:riemann-data} should reproduce the shock speed $s=1/2$ and the
linear profile $u=x/t$ inside the rarefaction fan, providing a
computational check on the exact formulas derived analytically above.
\end{remark}

\subsubsection{Numerical results}
\label{subsubsec:godunov-results}

\emph{Setup.} The scheme of
Sections~\ref{subsubsec:finite-volume}--\ref{subsubsec:godunov-flux} was
implemented (Python/NumPy; script \texttt{code/godunov.py}) on
$x\in[-2,2]$ with $N$ cells, initial cell averages of the Riemann data
(the discontinuity lies on a cell interface), zero-gradient (outflow)
boundary conditions, time step $\Delta t=0.45\,\Delta x/\max_j|U_j^n|$
(CFL number $0.45$ in \eqref{eq:cfl}, last step shortened to land on
$t=1$), and was run to $t=1$ on the rarefaction data
$(u_\ell,u_r)=(0,1)$ of Section~\ref{sec:nonuniqueness} and the shock data
$(u_\ell,u_r)=(1,0)$. The error is the $L^1$ error
$\|U-u\|_{L^1(-2,2)}=\int|U-u|\,dx$ against the exact entropy solution of
Theorem~\ref{thm:riemann-solution}, with the exact solution integrated
over each cell (so the cell containing the shock is not mis-scored).
Two sanity checks were passed: the case-based flux of
\eqref{eq:godunov-flux} agrees with the closed form
$\min_{[u_\ell,u_r]}f$ (if $u_\ell\le u_r$), $\max_{[u_r,u_\ell]}f$
(otherwise) on $10^5$ random state pairs, and total mass equals the initial
mass ($2$) plus the net boundary flux, namely $-0.5$ for the rarefaction
problem (final mass $1.5$) and $+0.5$ for the shock problem (final mass
$2.5$), to ten digits. Figure~\ref{fig:godunov-results} compares the numerical solution
with $N=400$ against the exact solution.

\begin{figure}[h]
\centering
\includegraphics[width=0.44\textwidth]{figures/godunov_rarefaction.pdf}
\includegraphics[width=0.44\textwidth]{figures/godunov_shock.pdf}
\caption{Godunov scheme (blue dots, $N=400$ cells) against the exact
entropy solution (red) of Theorem~\ref{thm:riemann-solution}, at
$t=1$. Left: rarefaction data $(u_\ell,u_r)=(0,1)$; the scheme
reproduces the linear fan $u=x/t$, with smearing at the two corners
(visibly larger at the upper corner $x=1$) typical of a first-order
scheme. Right: shock data $(u_\ell,u_r)=(1,0)$; the discontinuity is
located at the Rankine--Hugoniot position $x=st=0.5$
(eq.~\eqref{eq:rankine-hugoniot-burgers}) and is spread over about three
cells, with no overshoot or oscillation, as expected of a monotone
scheme.}
\label{fig:godunov-results}
\end{figure}

\emph{Convergence.} Table~\ref{tab:godunov-conv} and
Figure~\ref{fig:godunov-conv} report the $L^1$ error for
$N=200,400,\dots,6400$ (the choice of these resolutions is explained below).

\begin{table}[h]
\centering
\begin{tabular}{r r r r r r}
  $N$ & $\Delta x$ & $L^1$ (rarefaction) & order & $L^1$ (shock) & order \\\hline
  200  & $2.0\times10^{-2}$ & $3.05\times10^{-2}$ & --- & $9.91\times10^{-3}$ & --- \\
  400  & $1.0\times10^{-2}$ & $1.82\times10^{-2}$ & 0.74 & $4.98\times10^{-3}$ & 0.99 \\
  800  & $5.0\times10^{-3}$ & $1.07\times10^{-2}$ & 0.77 & $2.50\times10^{-3}$ & 0.99 \\
  1600 & $2.5\times10^{-3}$ & $6.14\times10^{-3}$ & 0.80 & $1.24\times10^{-3}$ & 1.01 \\
  3200 & $1.25\times10^{-3}$ & $3.48\times10^{-3}$ & 0.82 & $6.23\times10^{-4}$ & 0.99 \\
  6400 & $6.25\times10^{-4}$ & $1.95\times10^{-3}$ & 0.84 & $3.13\times10^{-4}$ & 0.99
\end{tabular}
\caption{$L^1$ errors of the Godunov scheme at $t=1$ and observed orders
$\log_2(e_{N/2}/e_N)$ between consecutive rows. All $N$ are multiples of
$8$, so the shock lies on a cell interface (see the text).}
\label{tab:godunov-conv}
\end{table}

\begin{figure}[h]
\centering
\includegraphics[width=0.44\textwidth]{figures/godunov_convergence.pdf}
\caption{$L^1$ error at $t=1$ against $\Delta x$ (log--log), with an
$O(\Delta x)$ reference line.}
\label{fig:godunov-conv}
\end{figure}

\emph{Shock: dependence on grid alignment.} The shock at $x=st=0.5$ lies
on a cell interface exactly when $0.5/\Delta x=N/8$ is an integer, i.e.\
when $8\mid N$. For such $N$ the shock error is $0.50\,\Delta x$ to two
digits (values $0.495$--$0.501$; the observed order is $\approx1.0$, as in
Table~\ref{tab:godunov-conv}). For other $N$ it depends on where the shock
sits within a cell: for $N=97,\dots,103$ the error ranges from
$0.60\,\Delta x$ to $0.77\,\Delta x$ (e.g.\ $0.655\,\Delta x$ at $N=100$,
$0.746\,\Delta x$ at $N=97$), against $0.50\,\Delta x$ at $N=96$ and
$N=104$. Observed orders between a misaligned and an aligned resolution are
therefore unreliable: $N=100\to200$ would give an order of $1.40$ for the
shock. This is why the table uses $N=200\cdot2^k$. In every case the error
is proportional to $\Delta x$, consistent with the discontinuity being
resolved over a fixed number of cells.

\emph{Rarefaction.} Here the error also decreases monotonically, but the
observed order is $\approx0.7$--$0.8$ and is still slowly increasing
($0.74,0.77,0.80,0.82,0.84$), so these experiments do not by themselves
establish a first-order rate. The larger error is attributable to the two
corners of the fan, where the exact solution is only Lipschitz and the
scheme's numerical diffusion smears it; we do not claim more than the data
show. (Theorem~\ref{thm:godunov-convergence} asserts convergence, not a
rate.) In both cases the computed shock position and fan slope agree with
the analytic values $s=1/2$ and $u=x/t$ of
Section~\ref{sec:riemann-problem}, which is the check proposed in
Remark~\ref{rem:godunov-check}.

\subsubsection{Numerical diffusion as vanishing viscosity}
\label{subsubsec:numerical-viscosity}

This subsection tests a heuristic that connects the scheme to the guiding
question: a monotone scheme converges to the \emph{entropy} solution because
its numerical diffusion acts as a vanishing viscosity (the mechanism of
Section~\ref{subsec:vanishing-viscosity}). It is a consistency check, not a
proof; the rigorous statement is Theorem~\ref{thm:godunov-convergence}.

In the shock problem all states lie in $[0,1]$ and $s=\tfrac12>0$, so the
Godunov flux is the upwind flux and the scheme reads $U_j^{n+1}=U_j^n-
\lambda(f(U_j^n)-f(U_{j-1}^n))$ with $\lambda=\Delta t/\Delta x$. A Taylor
expansion gives the modified equation
\begin{equation}
  \label{eq:modified-equation}
  u_t+u\,u_x=\frac{\Delta x}{2}\big(u(1-\lambda u)\,u_x\big)_x+O(\Delta x^2),
\end{equation}
a viscous Burgers' equation with state-dependent viscosity
$\nu(u)=\tfrac{\Delta x}2u(1-\lambda u)\ge0$ exactly when the CFL condition
\eqref{eq:cfl} holds. Its travelling shock from $1$ to $0$, with speed
$\tfrac12$, satisfies $\nu(u)\,u_\xi=\tfrac12u(u-1)$ in $\xi=x-t/2$, i.e.\ in
units of $\Delta x$
\begin{equation}
  \frac{du}{d\xi}=\frac{u-1}{1-\lambda u},\qquad
  \xi(u)=\xi_0+(1-\lambda)\ln(1-u)-\lambda u .
\end{equation}
The profile decays exponentially to $1$ as $\xi\to-\infty$ and reaches $0$ at
a finite $\xi_0$. Conservation of mass (the scheme is conservative) fixes
$\xi_0=1-\lambda/2$, and the $L^1$ distance from the sharp shock is
\begin{equation}
  \frac{\|u-u_{\rm exact}\|_{L^1}}{\Delta x}
  =2\int_0^{u_0}\frac{u(1-\lambda u)}{1-u}\,du,\qquad
  (1-\lambda)\ln(1-u_0)-\lambda u_0=-\Big(1-\frac\lambda2\Big).
\end{equation}
Table~\ref{tab:numerical-viscosity} compares this with the measured shock
error on aligned grids (here $\lambda$ equals the CFL number, since
$\max|U|=1$); the code is \texttt{code/numerical\_viscosity.py}.

\begin{table}[h]
\centering
\begin{tabular}{c c c c}
  CFL $=\lambda$ & measured ($N=1600$) & predicted & difference \\\hline
  0.10 & 0.627 & 0.686 & $+9.4\%$ \\
  0.20 & 0.589 & 0.635 & $+8.0\%$ \\
  0.30 & 0.553 & 0.586 & $+5.8\%$ \\
  0.45 & 0.496 & 0.511 & $+3.1\%$ \\
  0.60 & 0.448 & 0.438 & $-2.4\%$ \\
  0.80 & 0.352 & 0.342 & $-3.0\%$ \\
  0.95 & 0.300 & 0.272 & $-9.3\%$
\end{tabular}
\caption{Shock $L^1$ error divided by $\Delta x$: measured Godunov (grids with
$8\mid N$) against the modified-equation prediction.}
\label{tab:numerical-viscosity}
\end{table}

The prediction is within $3\%$ at the CFL number $0.45$ used above and within
$10\%$ across the whole range, and it reproduces the observed decrease of the
error with $\lambda$. At CFL $0.95$ the measured value is still resolution
dependent ($0.334$ at $N=400$, $0.300$ at $N=1600$), so that row is not
asymptotic. Equating the error to the viscous-shock value $4\ln2\,\nu_{\rm eff}$
of Proposition~\ref{prop:viscous-profile} (valid for any shock speed, here
$s=\tfrac12$, $\Delta=1$, by Galilean invariance) defines an effective viscosity: the
measured $0.50\,\Delta x$ gives $\nu_{\rm eff}\approx0.180\,\Delta x$, the
modified-equation profile gives $0.184\,\Delta x$, and the cruder estimate
$\tfrac{\Delta x}2u(1-\lambda u)$ at $u=\tfrac12$ gives $0.194\,\Delta x$.
Since $\nu_{\rm eff}=O(\Delta x)\to0$, in the shock layer the scheme behaves
like a vanishing viscosity, which is the numerical counterpart of
Theorem~\ref{thm:vanishing-viscosity}. This agreement should not be over-read:
the numerical shock occupies only a few cells, and
\eqref{eq:modified-equation} neglects the $O(\Delta x^2)$ terms.